\scanpage{9}
\uline{Schémas quasi-affines} \inserted{(\unsure{quasicompacts|quasi-compacts}, et isom. à sous-schéma ouvert d'un schéma affine)} / \uline{morphismes quasi-affines}

Conditions équivalentes sur schéma $X$ cohérent, $A=\Gamma(X,\underline O_X)$

a) $X$ quasi-affine

b) $X\overset{u}{\mathrel{-}}\operatorname{Spec}A$ une imm. ouverte

b') $X\mathrel{-}\operatorname{Spec}A$ un homéomorphisme dans

c) $\underline O_X$ est ample

\[
\left\{\begin{aligned}
 &\text{c') Les }X_f\text{ forment une base top. de }X\\
 &\text{c'') les }X_f\text{ affines forment base top. de }X\\
 &\text{c''') Tout } F\text{ quasi-cohérent \uline{engendré} par ses sections}\\
 &c^{\mathrm{iv})})\ \text{Id pour }F\text{ qu. ch }\subset\underline O_X
\end{aligned}\right.
\]

d) $\exists$ Faisceau \deleted{ample} \inserted{inv.} $L$ tel que $L$ et $L^{-1}$ amples.

d') Tout Module inv. $L$ est ample.

\uline{Morphismes} \uline{quasi-affines} : Stable par changemt de base, par composition les imm. quasi-compactes sont qu. affines, donc si $g$ qu. sép, $\delta_g$ qu. affine

$gf$ quasi-affine et $g$ qu. sép $\Longrightarrow f$ quasi-affine

\uline{Th} Si $f$ de type fini, \deleted{\unsure{séparé}} à fibres finies ($f$ quasi-fini) et $f$ séparé, donc $f$ quasi-affine (donc quasi-projectif)

\uline{morph. quasi-projectif} : de t.f et $\exists L$ $f$-ample.

Propre $={}$ univers. fermé, séparé, de t.f.

entier

\[
\begin{aligned}\text{projectif}&=\text{propre}+\text{quasi-projectif}\\
 &=\deleted{\unsure{\text{propre}}}\ \text{fermé}+\text{quasi-projectif}
\end{aligned}
\]

\diagram{\begin{tikzcd}[ampersand replacement=\&,column sep=.45em,row sep=1.05em,cells={font=\scriptsize},execute at end picture={\draw ([xshift=-.45em,yshift=.55em]\tikzcdmatrixname-1-2.north west) rectangle ([xshift=.45em,yshift=-.5em]\tikzcdmatrixname-3-7.south east);}]
 {} \&\text{imm. fermée}\arrow[rr,Rightarrow]\arrow[d,Rightarrow] \&{} \&\text{fini}\arrow[rr,Rightarrow]\arrow[dl,Rightarrow]\arrow[d,Rightarrow] \&{} \&\text{entier}\arrow[r,Rightarrow]\arrow[d,Rightarrow] \&\text{affine}\arrow[dd,Rightarrow]\\
 {} \&\text{imm. qu. cpt}\arrow[d,Rightarrow] \&\substack{\text{affine}\\\text{de t.f.}}\arrow[dl,Rightarrow] \&\text{proj}\arrow[r,Rightarrow]\arrow[d,Rightarrow] \&\text{propre}\arrow[r,Rightarrow]\arrow[drr,Rightarrow] \&\text{univ. fermé} \&{}\\
 \deleted{\substack{\text{affine}\\\text{de t.f.}}}\arrow[r,Rightarrow]\&\substack{\text{quasi-affine}\\\text{de t.f.}}\arrow[rr,Rightarrow] \&{} \&\text{qu. proj}\arrow[rrr,Rightarrow] \&{} \&{} \&\text{séparé}
\end{tikzcd}}

\uline{Th.} \textcircled{a} \uline{fini} $={}$ \uline{propre à fibres finies} (\unsure{ensemblistement|cas noethérien}) $={}$ projectif à fibres finies

\deleted{\textcircled{b} \uline{séparé} \unsure{quasi-fini|quasi-} $\Rightarrow$ \uline{\unsure{qu. affine|affine}} \unsure{à} \uline{fibres finies}}

% archive sheet number: 8
